\graphicspath{{abstracts/workshops/02-installation-konfiguration/img/}{img/}}

\begin{beitrag}{WissKI installieren und konfigurieren: Ein Hands-on-Workshop für Projektstarter*innen}{%
Petra Pohl\,\orcidlink{0000-0001-7890-1234} \\
\small \href{mailto:petra.pohl@example.org}{petra.pohl@example.org} \\
\small Universität Leipzig \\[0.5cm]
Bernd Braun\,\orcidlink{0000-0002-3456-7890} \\
\small \href{mailto:bernd.braun@example.org}{bernd.braun@example.org} \\
\small Friedrich-Alexander-Universität Erlangen-Nürnberg
}

\begin{abstract}
Der Hands-on-Workshop führt Projektstarter*innen in fünf Blöcken durch alle Kernschritte einer WissKI-Installation: Server-Vorbereitung, Drupal- und WissKI-Modulinstallation (Composer, Drush), Triplestore-Anbindung (GraphDB oder Fuseki), Pathbuilder-Basiskonfiguration sowie Backup-, Monitoring- und Update-Routinen. Alle Konfigurationsschritte sind reproduzierbar dokumentiert; die Workshop-Materialien sind öffentlich verfügbar, und die Teilnehmer*innen erhalten für drei Monate einen Support-Kanal. Am Ende verfügen sie über eine lauffähige, triplestore-verbundene WissKI-Installation mit dem Verständnis für jeden einzelnen Installationsschritt.
\end{abstract}

\section{Zielsetzung}
Der Workshop richtet sich an Personen, die in den kommenden Monaten ein WissKI-Projekt starten oder eine bestehende Installation in einen produktiven Betrieb überführen wollen. Im Zentrum stehen die Installation einer aktuellen WissKI-Version, die Grundkonfiguration auf Server-Ebene, das Aufsetzen eines verbundenen Triplestores und die ersten Schritte der inhaltlichen Konfiguration. Wir verfolgen einen \emph{Hands-on}-Ansatz: Jede*r Teilnehmer*in arbeitet mit der eigenen Installation und schließt den Workshop mit einer betriebsbereiten, dokumentierten WissKI-Umgebung ab.

\section{Voraussetzungen}
Wir setzen grundlegende Kenntnisse der Linux-Kommandozeile (\texttt{ssh}, einfache Paketverwaltung, \texttt{systemctl}) und ein Verständnis der gängigen Webdienst-Begriffe (Reverse-Proxy, virtuelle Hosts, TLS) voraus. Eine vorbereitete virtuelle Maschine mit Debian 12 und einem zweiten Container für den Triplestore wird bereitgestellt; alternativ können Teilnehmer*innen mit administrativem Zugriff auf eine eigene Linux-Umgebung arbeiten. Eine schrittweise Anleitung wird zwei Wochen vor dem Termin versendet.

\section{Ablauf des Workshops}
Der Workshop ist in fünf zeitliche Blöcke von je etwa eine Stunde gegliedert.

Block eins befasst sich mit der Vorbereitung der Server-Umgebung: Systempaket-Aktualisierungen, Anlegen der Dienste-Nutzer, Einrichtung eines Reverse-Proxys, Konfiguration der Datenbank-Komponenten. Wir legen besonderen Wert auf eine saubere Trennung der Konfigurationsdateien, sodass spätere Änderungen reproduzierbar nachvollzogen werden können.

Block zwei behandelt die eigentliche Installation von Drupal und der WissKI-Module über Composer und Drush. Diese Phase ist die technisch heikelste; wir werden auf die in der einschlägigen Literatur dokumentierten Stolpersteine im Detail eingehen.\footcite[48]{pohl2024-installation-leitfaden}

Block drei widmet sich der Anbindung an den Triplestore. Wir installieren GraphDB (alternativ Apache Jena Fuseki), legen ein erstes Repositorium an und konfigurieren die WissKI-Adapter so, dass Schreib- und Lesepfade zuverlässig funktionieren. Die Konfigurationsentscheidungen werden ausführlich besprochen, weil sie in der Betriebsphase nur mit erheblichem Aufwand revidiert werden können.

Block vier beschäftigt sich mit der Basiskonfiguration der WissKI-Anwendung: Anlegen der Inhalts-Bundles, Konfiguration der Pathbuilder, Anlegen der ersten Formulare. Wir folgen dabei einer Konfigurationsvorlage, die im WissKI-Konsortium für neu startende Projekte abgestimmt ist und die einen guten Kompromiss zwischen unmittelbarer Nutzbarkeit und langfristiger Anpassbarkeit darstellt.\footcite[vgl.][]{braun2024-konfigurationsvorlage}

Block fünf widmet sich Inbetriebnahme-Aspekten, die häufig zu spät bedacht werden: Backup-Strategien, Monitoring, Update-Routinen. Wir geben für jede dieser Aspekte eine konkrete Empfehlung, deren Umsetzung wir im Workshop einmal exemplarisch durchspielen. Die Empfehlungen sind so gewählt, dass sie auch mit überschaubaren Personalressourcen dauerhaft tragfähig sind.

\section{Materialien und Nachbetreuung}
Sämtliche im Workshop verwendeten Skripte, Konfigurationsdateien und Anleitungen werden in einem öffentlichen Git-Repositorium bereitgestellt. Die Teilnehmer*innen erhalten darüber hinaus für die ersten drei Monate nach dem Workshop einen niedrigschwelligen Support-Kanal -- per Mailingliste und in Form eines monatlichen offenen Online-Treffens, das mit konkreten Fragen aus der Inbetriebnahme-Phase besucht werden kann.

\section{Erwartetes Ergebnis}
Am Ende des Workshops verfügen die Teilnehmer*innen über eine eigene, lauffähige WissKI-Installation, die mit einem produktiven Triplestore verbunden ist und die Grundlagen einer projektspezifischen Konfiguration bereits enthält. Wichtiger noch: Sie haben einmal jeden Schritt der Installation selbst nachvollzogen und sind dadurch in der Lage, im Betrieb gezielt zu reagieren, wenn etwas nicht wie erwartet funktioniert.

\end{beitrag}

\section*{Kurzbiografie(n)}

\subsection*{Petra Pohl}
Petra Pohl ist Systemadministratorin und IT-Beraterin am Universitätsrechenzentrum Leipzig. Sie betreut die produktiven Drupal- und WissKI-Installationen der Leipziger Geisteswissenschaften und hat in dieser Funktion zahlreiche Neuinstallationen begleitet. Ihre Schwerpunkte liegen in der Inbetriebnahme nachhaltiger Server-Konfigurationen und in der Vermittlung des dafür benötigten Wissens.

\subsection*{Bernd Braun}
Bernd Braun ist technischer Koordinator des WissKI-Konsortiums an der FAU Erlangen-Nürnberg. Seine Schwerpunkte liegen in der Weiterentwicklung des WissKI-Pathbuilders, in der Konfiguration und Inbetriebnahme produktiver Installationen sowie in der Unterstützung neu startender Projekte. Er war an der Entwicklung der aktuellen WissKI-Konfigurationsvorlage maßgeblich beteiligt und engagiert sich in der Open-Source-Pflege des WissKI-Quellcodes.
